\documentclass[10pt,a4paper]{amsart}
\usepackage[margin=19mm]{geometry}
\usepackage{amsmath,amssymb,mathtools}
\usepackage[scr=rsfs,cal=euler]{mathalfa}
\usepackage{xfrac}
\usepackage[unicode,hidelinks,bookmarks=false]{hyperref}
\newcommand{\ie}{i.e.\ }
\newcommand{\cf}{cf.\ }
\newcommand{\resp}{respectively\ }
\newcommand{\Subsection}{\subsection}
\providecommand{\nobreakdash}{\nobreak\hbox{-}}
\setlength{\emergencystretch}{2em}
\multlinegap0pt
\usepackage{bm}
\usepackage{bbm}
\usepackage{dsfont}
\usepackage{xspace}
\usepackage{cancel}
\usepackage{mathtools}
\usepackage{amsthm}
\makeatletter
\@ifundefined{theorem}{\newtheorem{theorem}{Theorem}}{}
\@ifundefined{lemma}{\newtheorem{lemma}[theorem]{Lemma}}{}
\makeatother
\DeclareMathOperator{\Mat}{Mat}
\DeclareMathOperator{\ord}{ord}
\newcommand{\equivk}{\overset{\mathcal{K}}{\sim}}
\newcommand{\naxV}{\mathfrak{n}_V}
\title[Detecting and Quantifying Isolated Singularities over Discre]{Detecting and Quantifying Isolated Singularities over Discrete Valuation Rings}
\author{Yotam Svoray}
\date{Source: arXiv:2512.18506v2 (CC BY 4.0)}
\begin{document}
\maketitle

\noindent\emph{Source and licence.} Detecting and Quantifying Isolated Singularities over Discrete Valuation Rings. Version arXiv:2512.18506v2 (CC BY 4.0), distributed under the Creative Commons Attribution 4.0 International licence, as recorded in the arXiv licence metadata for this exact version and shown on the abstract page https://arxiv.org/abs/2512.18506v2. Full rights notice: licenses/arxiv-2512.18506-CC-BY-4.0.txt.\par\smallskip

\noindent\textit{Editorial scope.} Counted unit: the Lemma whose source label is \texttt{prop:rank\_good\_def} in the article (source0.tex lines 1129--1131, source label \texttt{prop:rank\_good\_def}), delivered together with the article's own complete original proof of it (source0.tex lines 1133--1170, one \texttt{proof} environment of 38 lines). No mathematics is written, repaired or re-derived: statement and proof are the article's own text line for line. Required context retained uncounted: lem:W[[x]]\_good\_prop (lines 397--413). Every definition, display and label of those lines is retained unchanged. Omitted from this package and not redistributed: 1--396, 414--1128, 1132--1132, 1171--1609. Nothing inside a retained excerpt is omitted or altered: every retained body is delivered exactly as the article's lines are written, so no omission receipt of any kind accompanies this package. Citations: the retained excerpts carry no citation command at all, so no bibliography entry of the article is transcribed and no reference is redistributed. Reference adaptation: none was needed and none was made; every reference inside the delivered unit resolves inside it, so no unresolved reference or citation is produced. Printed slips kept literal: no printed word, symbol, hypothesis or number of the counted statement or of its proof was altered. Layout and class: the article's own class is \texttt{article} with packages including inputenc, authblk, hyperref, enumitem, array, tikz, tikz-cd, amsthm, amsmath, amssymb, amsfonts, faktor, fontenc, xfrac; the wrapper is the lane's pinned \texttt{amsart} editorial wrapper at 10pt on A4, which restates the theorem-like environments the retained text needs and adds the article's own macro definitions that the retained text needs (\texttt{\textbackslash Mat}, \texttt{\textbackslash equivk}, \texttt{\textbackslash naxV}, \texttt{\textbackslash ord}), with extra wrapper packages (bm, bbm, dsfont, xspace, cancel, mathtools). The delivered numbering is the unit's own. Assets: no retained span contains a figure environment, a graphics-inclusion command, a PDF-page inclusion, a table or a TikZ picture; no image file is copied into this package, the package declares no asset dependency and no image file is redistributed with it. Licence: version arXiv:2512.18506v2 of this article is distributed under the Creative Commons Attribution 4.0 International licence, as recorded in the arXiv licence metadata for this exact version and shown on the abstract page https://arxiv.org/abs/2512.18506v2. A full rights notice is delivered at licenses/arxiv-2512.18506-CC-BY-4.0.txt. Characters: every retained excerpt is pure ASCII, so the delivered engine, XeLaTeX with its default Latin Modern text font, reports no missing character.
\begin{lemma}\label{lem:W[[x]]_good_prop}
\begin{enumerate}
    \item \textup{(Implicit Function Theorem)}: Given a sequence of power series $f_1, \dots, f_m \in R[[\underline{x},y_1, \dots y_m]]$ such that $f_i(\underline{0})=0$ for every $i$ and 
    \begin{equation*}
    \det \left(
        \begin{bmatrix}
            \frac{\partial f_1}{\partial y_1}(\underline{0}) & \cdots & \frac{\partial f_1}{\partial y_m}(\underline{0})\\
            \vdots & & \vdots \\
            \frac{\partial f_m}{\partial y_1}(\underline{0}) & \cdots & \frac{\partial f_m}{\partial y_m}(\underline{0})
        \end{bmatrix}
        \right) \notin \mathfrak{m},
    \end{equation*}
    then we can find a sequence of power series $g_1, \dots, g_m \in \mathfrak{M} \subset R[[\underline{x}]]$ for which we have that $f_i\left(x_1, \dots, x_n, g_1, \dots g_m\right)=0$ for every $i$. 
    \item \textup{(Inverse Function Theorem)}: Let $f_1, \dots, f_n \in \mathfrak{M} \subset R[[\underline{x}]]$. Then the  Jacobian matrix $Df=\left[\frac{\partial f_i}{\partial x_j}\left(0\right)\right]_{i,j} \in \Mat_n\left(R\right)$ is an invertible matrix if and only if the unique $R-$homomorphism $\varphi \colon R[[\underline{x}]] \to R[[\underline{x}]]$ defined by $\varphi\left(x_i\right)=f_i$ defines an $R-$automorphism. 
    %\item (\textbf{Approximation Theorem}): Let $\left(f_1, \dots, f_N\right) \in \Vx^N$ and let $\{Y_c\}_{c=1}^\infty \subset \Vx^N$. Assume that there exists a monotonically increasing function $\beta \colon \mathbb{N} \to \mathbb{N}$ such that $f_i\left(Y_c\right) \in \maxV^{\beta\left(c\right)}$ for every $i$ and for every $c$. Then, there exists some $Y \in \Vx^N$ such that $Y - Y_c \in \maxV^c$ and $f_i\left(Y\right)=0$ for every $i$. 
\end{enumerate}
\end{lemma}

\begin{lemma}\label{prop:rank_good_def}
    Let $f, g \in \mathfrak{a}_3 \subset V[[x_1, \dots, x_l,y]]$. Then we have  $f\left(x_1, \dots, x_l\right) \equivk  g\left(x_1, \dots, x_l\right)$ in $V[[x_1, \dots, x_l]]$ if and only if in $V[[\underline{x},y]]$ we have the equivalence $$f\left(x_1, \dots, x_l\right)+y^2+x_{l+1}^2 + \dots + x_n^2 \equivk  g\left(x_1, \dots, x_l\right)+y^2+x_{l+1}^2 + \dots + x_n^2.$$  
\end{lemma}

\begin{proof}
    First, if $f\left(x_1, \dots, x_l\right) \equivk g\left(x_1, \dots, x_l \right)$ then we can find some automorphism $\varphi$ of $V[[x_1, \dots, x_l, y]]$ and some unit $u \in V[[x_1, \dots, x_l, y]]$ such that $f = u \varphi(g)$. So, by extending $\varphi$ to an automorphism $\psi$ of $V[[\underline{x},y]]$ by setting $\psi \left(x_i\right)=v^{-1} x_i$ for every $i>l$, where $v^2=u$, we can conclude that 
    \begin{equation*}
        f\left(x_1, \dots, x_l\right) +x_{l+1}^2 + \dots + x_n^2  = u \psi\left(g\left(x_1, \dots, x_l\right) +x_{l+1}^2 + \dots + x_n^2\right). 
    \end{equation*}
    

    \noindent Second, assume that 
    \begin{equation*}
        F=f\left(x_1, \dots, x_l \right)+ y^2+ x_{l+1}^2 + \dots + x_n^2 \equivk g\left(x_1, \dots, x_l\right)+y^2 +x_{l+1}^2 + \dots + x_n^2=G.
    \end{equation*}
    Therefore, there exists some $\varphi_1, \dots, \varphi_{n+1} \in V[[\underline{x},y]]$ and some unit $u$ such that 
    \begin{equation*}
         g\left(\varphi_1, \dots, \varphi_{l}\right) + \sum_{i=l+1}^{n+1} \varphi_i^2 =   uf\left(x_1, \dots, x_l\right) + u\cdot \sum_{i=l+1}^n x_i^2 +uy^2.
    \end{equation*}
    
    \noindent Now,  write $\psi_i = u^{-0.5} \varphi_i = a_i + h_i$ where $a_i$ is a linear term over $\langle x_1, \dots, x_n,y \rangle$ and $h_i \in \naxV^2$. Then  we have that $$\sum_{i=l+1}^{n+1} a_i^2 - \sum_{i=l+1}^{n+1} x_i^2 - y^2 \in \naxV^2.$$ In addition, we have that 
    \begin{equation*}
        \sum_{i=l+1}^{n+1} \psi_i^2 = \sum_{i=l+1}^{n+1} a_i^2 + \sum_{i=l+1}^{n+1} h_i(2a_i+h_i),
    \end{equation*}
    and since $\ord_{\langle x_1, \dots, x_n,y \rangle}(a_{l+1}) =  \cdots =\ord_{\langle x_1, \dots, x_n,y \rangle}(a_{n+1})=1$,  we can conclude that $\sum_{i=1}^{n+1} a_i^2 = \sum_{i=l+1}^{n} x_i^2 +y^2$. Now, if we look at $F_i = 2a_i + h_i $ (recalling that the characteristic of $\kappa$ is not $2$) then  
    %
    \begin{equation*}
        \det\left( \left[\frac{\partial F_i}{\partial x_j} \left(0\right)\right]_{n+1 \geq i,j \geq l+1}\right) =2\det\left( \left[\frac{\partial a_i}{\partial x_j} \left(0\right)\right]_{n+1 \geq i,j \geq l+1}\right) \in V^{\times},
    \end{equation*}
    %
    \noindent where $x_{n+1}=y$. Therefore, by the implicit function theorem and the inverse function theorem (Lemma~\ref{lem:W[[x]]_good_prop}), there exists some $\theta_{l+1}, \dots, \theta_{n+1} \in V[[x_1, \dots, x_l]]$ for which we have that $F_i\left(x_1, \dots, x_l, \theta_{l+1}, \dots, \theta_{n+1}\right)=0$ for every $i$. For every $i \leq l$ define  $\beta_i = \psi_i\left(x_1, \dots, x_l, \theta_{l+1}, \dots, \theta_{n+1}\right)$. Then the map $\beta$ on $V[[x_1, \dots, x_l,y]]$ defined by $\beta\left(x_i\right)=\beta_i$ must be an automorphism. This tells us that $\left(\psi_1, \dots, \psi_l, a_{l+1}+\frac{h_{l+1}}{2}, \dots,  a_{n+1}+\frac{h_{n+1}}{2} \right)$ defines a $V-$automoprhism of $V[[\underline{x},y]]$. Thus, by the inverse function theorem (Lemma~\ref{lem:W[[x]]_good_prop}) we have that $\langle \psi_1, \dots, \psi_l, a_{l+1}+\frac{h_{l+1}}{2}, \dots,  a_{n+1}+\frac{h_{n+1}}{2} \rangle = \naxV$. Therefore, for every $i \leq l$ there exists some $\alpha_{1,i}, \dots,\alpha_{n+1,i}$ such that 
    \begin{equation*}
        x_i = \sum_{j=1}^l \alpha_{j,i} \psi_j + \sum_{j=l+1}^{n+1} \alpha_{j,i} \left(a_{j}+\frac{h_{j}}{2}\right), 
    \end{equation*}
    which gives us that 
    %
    \begin{equation*}
        x_i=\sum_{j=1}^l \alpha_j\left(x_1, \dots, x_l, \theta_{l+1}, \dots, \theta_{n+1}\right) \beta_i. 
    \end{equation*}
    %
    Therefore we have that $\langle \beta_1, \dots, \beta_l \rangle = \langle x_1, \dots, x_l \rangle$, and by the inverse function theorem (Lemma~\ref{lem:W[[x]]_good_prop}), we can conclude that $\beta(f)=f\left(\beta_1, \dots, \beta_l\right)= u\left(g\left(x_1, \dots, x_l\right)\right)$ and the result follows.
\end{proof}

\end{document}
